\documentclass[11pt,a4paper]{article}

\usepackage[utf8]{inputenc}
\usepackage[T1]{fontenc}
\usepackage[turkish]{babel}
\usepackage[a4paper,margin=2.4cm]{geometry}
\usepackage{xcolor}
\usepackage{enumitem}
\usepackage{hyperref}
\usepackage{listings}
\usepackage{tcolorbox}
\usepackage{titling}
\usepackage{titlesec}
\usepackage{parskip}

% Renkler
\definecolor{brand}{HTML}{1E3A8A}
\definecolor{brandAcc}{HTML}{F97316}
\definecolor{soft}{HTML}{F8FAFC}
\definecolor{warn}{HTML}{B45309}
\definecolor{ok}{HTML}{047857}

\hypersetup{
  colorlinks=true,
  linkcolor=brand,
  urlcolor=brand,
  pdftitle={Alfa Emlak — 101evler ve hangiev XML Entegrasyon Rehberi},
  pdfauthor={Alfa Emlak}
}

\titleformat{\section}{\Large\bfseries\color{brand}}{\thesection.}{0.6em}{}
\titleformat{\subsection}{\large\bfseries\color{brand}}{\thesubsection.}{0.6em}{}

\lstdefinestyle{xmlstyle}{
  basicstyle=\ttfamily\footnotesize,
  backgroundcolor=\color{soft},
  frame=single,
  rulecolor=\color{black!10},
  framerule=0.4pt,
  breaklines=true,
  showstringspaces=false,
  literate={ş}{{\c s}}1 {ğ}{{\u g}}1 {ı}{{\i}}1 {İ}{{\.I}}1 {Ş}{{\c S}}1 {Ğ}{{\u G}}1 {ç}{{\c c}}1 {Ç}{{\c C}}1 {ö}{{\"o}}1 {Ö}{{\"O}}1 {ü}{{\"u}}1 {Ü}{{\"U}}1
}

\newtcolorbox{infobox}[1][]{
  colback=soft, colframe=brand, boxrule=0.5pt, arc=2pt,
  left=8pt, right=8pt, top=6pt, bottom=6pt,
  fonttitle=\bfseries\color{white},
  coltitle=white, colbacktitle=brand,
  title=#1
}

\newtcolorbox{warnbox}[1][]{
  colback=yellow!8, colframe=warn, boxrule=0.5pt, arc=2pt,
  left=8pt, right=8pt, top=6pt, bottom=6pt,
  title=#1, coltitle=white, colbacktitle=warn,
  fonttitle=\bfseries
}

\newtcolorbox{okbox}[1][]{
  colback=green!5, colframe=ok, boxrule=0.5pt, arc=2pt,
  left=8pt, right=8pt, top=6pt, bottom=6pt,
  title=#1, coltitle=white, colbacktitle=ok,
  fonttitle=\bfseries
}

\title{\vspace{-1cm}\Huge\color{brand}\textbf{Alfa Emlak}\\[4pt]\Large\color{black}\normalfont XML İlan Aktarımı — Entegrasyon Rehberi}
\author{\normalfont\large 101evler ve hangiev portalları için teknik dokümantasyon}
\date{\today}

\begin{document}

\maketitle
\thispagestyle{empty}

\vspace{-0.5cm}

\begin{infobox}[Site Bilgisi]
\textbf{Site adresi:} \url{https://alfa-emlak-websitesi.onrender.com/tr} \\
\textbf{Doküman amacı:} Alfa Emlak portföyündeki ilanların 101evler ve hangiev platformlarına XML feed üzerinden otomatik aktarımının nasıl yapılacağını anlatmaktır. \\
\textbf{Hedef kitle:} 101evler ve hangiev teknik / entegrasyon ekipleri.
\end{infobox}

\tableofcontents
\vspace{0.5cm}
\hrule
\vspace{0.5cm}

% ============================================================
\section{Genel Bakış}

Alfa Emlak; web sitesinde yayında olan tüm ilanları, her iki portala özel XML uç noktaları (\textbf{feed endpoint}) üzerinden \textbf{anlık} olarak sunar. Portal tarafının yapması gereken yalnızca aşağıda belirtilen URL'leri belirli bir aralıkla (önerilen: 15--30 dakika) çağırıp dönen XML çıktısını işlemektir. Veri tabanı erişimi, FTP, e-posta veya manuel yükleme gerekmez.

\begin{itemize}[leftmargin=*]
\item \textbf{Format:} XML (UTF-8 kodlamalı)
\item \textbf{Yöntem:} HTTPS GET
\item \textbf{Kimlik doğrulama:} Sorgu parametresi olarak verilen statik \textbf{token}
\item \textbf{Önbellek:} 5 dakika (sunucu tarafı). Bu süre boyunca cevap aynı kalır.
\item \textbf{Diller:} Başlık ve açıklamalar TR/EN/RU/DE/FA olarak çoklu dilde gönderilir.
\end{itemize}

% ============================================================
\section{Feed Adresleri}

\subsection{101evler için}

\begin{lstlisting}[style=xmlstyle]
GET https://alfa-emlak-websitesi.onrender.com/api/feeds/101evler.xml?token=<TOKEN>
\end{lstlisting}

\subsection{hangiev için}

\begin{lstlisting}[style=xmlstyle]
GET https://alfa-emlak-websitesi.onrender.com/api/feeds/hangiev.xml?token=<TOKEN>
\end{lstlisting}

\begin{warnbox}[TOKEN değerleri]
Her iki uç nokta için tarafınıza \textbf{özel olarak iletilecek} statik bir token kullanılır. Bu token gizli tutulmalıdır; üçüncü kişilerle paylaşılmamalıdır. Talep üzerine değiştirilebilir.

\smallskip
\textbf{Yanlış / eksik token:} \texttt{HTTP 401 Unauthorized} \\
\textbf{Token sunucuda tanımlı değil:} \texttt{HTTP 500} (operasyon tarafına bildirin)
\end{warnbox}

% ============================================================
\section{İstemci Tarafının Yapması Gerekenler}

\begin{enumerate}[leftmargin=*]
\item Yukarıdaki URL'ye, tarafınıza özel \texttt{token} parametresi ile HTTPS GET isteği atın.
\item Cevap \texttt{Content-Type: application/xml; charset=utf-8} olarak gelir.
\item Kök etiket \textbf{101evler için} \texttt{<ads>}, \textbf{hangiev için} \texttt{<listings>}'tir. Her ilan kök etiketin altında ayrı bir \texttt{<ad>} (101evler) veya \texttt{<listing>} (hangiev) çocuğu olarak bulunur.
\item İlanın benzersiz kimliği olarak \texttt{ad\_key} / \texttt{listing\_id} alanını kullanın. Aynı kimlikle tekrar gelen bir ilan \textbf{güncelleme} olarak değerlendirilmelidir; yeni kayıt açılmamalıdır.
\item Bir önceki feed çekiminizde gelmiş ancak güncel çekimde \textbf{bulunmayan} ilanlar; sitemizde yayından kaldırılmış veya silinmiş demektir. Portal tarafında bunları da kaldırmanız beklenir.
\item Önerilen çağrı sıklığı: \textbf{15 ila 30 dakikada bir}. Daha sık çağrılar, sunucu önbelleği nedeniyle aynı içeriği döndürür ve avantaj sağlamaz.
\end{enumerate}

\begin{okbox}[Tavsiye — Hata yönetimi]
HTTP 200 dışı yanıtlarda mevcut portföyü silmeyin; bir sonraki başarılı çekimi bekleyin. Bu yaklaşım, geçici ağ/servis problemlerinde ilan kaybını önler.
\end{okbox}

% ============================================================
\section{XML Yapısı — Genel Çerçeve}

\subsection{101evler örnek çıktı}

\begin{lstlisting}[style=xmlstyle]
<?xml version="1.0" encoding="UTF-8"?>
<ads>
  <ad>
    <lastupdate>2026-05-12 14:21:00</lastupdate>
    <type_id>...</type_id>
    <first_realtor_id>...</first_realtor_id>
    <second_realtor_id>...</second_realtor_id>
    <area_id>...</area_id>
    <reference_no>AE-2026-12345</reference_no>
    <title_type_id>...</title_type_id>
    <sale_or_rent>S</sale_or_rent>     <!-- S: Satilik, R: Kiralik -->
    <ad_title_tr>...</ad_title_tr>
    <ad_title_en>...</ad_title_en>
    <ad_description_tr>...</ad_description_tr>
    <price>185000</price>
    <price_for>T</price_for>           <!-- T: Toplam, M: Aylik vb. -->
    <currency>602</currency>           <!-- 601=TRY, 602=USD, 603=EUR, 604=GBP -->
    <room_count_id>...</room_count_id>
    <is_new>N</is_new>
    <bedroom_count>3</bedroom_count>
    <bathroom_count>2</bathroom_count>
    <floor>2</floor>
    <total_area>120</total_area>
    <build_age_id>...</build_age_id>
    <furnishing_id>...</furnishing_id>
    <map_available>1</map_available>
    <lat>35.1856</lat>
    <lng>33.3823</lng>
    <ad_key>123e4567-e89b-12d3-a456-426614174000</ad_key>
    <url>https://alfa-emlak-websitesi.onrender.com/tr/ilan/AE-2026-12345</url>
    <ad_specs>...</ad_specs>
    <ad_pictures>
      <ad_picture>
        <url>https://.../image1.jpg</url>
        <sort_order>0</sort_order>
      </ad_picture>
      ...
    </ad_pictures>
  </ad>
</ads>
\end{lstlisting}

\subsection{hangiev örnek çıktı}

\begin{lstlisting}[style=xmlstyle]
<?xml version="1.0" encoding="UTF-8"?>
<listings>
  <listing>
    <portal_id>...</portal_id>
    <agent_id>...</agent_id>
    <office_id>...</office_id>
    <listing_id>AE-2026-12345</listing_id>
    <last_update>2026-05-12 14:21:00</last_update>
    <listing_status>active</listing_status>
    <sale_or_rent>sale</sale_or_rent>  <!-- sale | rent -->
    <title_tr>...</title_tr>
    <description_tr>...</description_tr>
    <price>185000</price>
    <currency>EUR</currency>           <!-- TRY | USD | EUR | GBP -->
    <property_type>apartment</property_type>
    <rooms>3+1</rooms>
    <area>120</area>
    <bedrooms>3</bedrooms>
    <bathrooms>2</bathrooms>
    <floor>2</floor>
    <lat>35.1856</lat>
    <lng>33.3823</lng>
    <url>https://alfa-emlak-websitesi.onrender.com/tr/ilan/AE-2026-12345</url>
    <photos>
      <photo>https://.../image1.jpg</photo>
      ...
    </photos>
  </listing>
</listings>
\end{lstlisting}

\begin{infobox}[Tam alan listesi]
Yukarıdaki örnekler kısaltılmıştır. Üretilen XML; tüm fotoğraflar, özellik etiketleri (\texttt{ad\_specs} / \texttt{features}), referans numarası, koordinatlar, çoklu dil başlık ve açıklamalarını içerir. \textbf{Lütfen entegrasyona başlamadan önce gerçek bir çıktı örneği talep edin}; tarafınıza canlı XML örneği iletilecektir.
\end{infobox}

% ============================================================
\section{Desteklenen Para Birimleri}

\begin{tabular}{lll}
\textbf{Para birimi} & \textbf{101evler kodu} & \textbf{hangiev kodu} \\
Türk Lirası & 601 & TRY \\
ABD Doları & 602 & USD \\
Avro & 603 & EUR \\
İngiliz Sterlini & 604 & GBP \\
\end{tabular}

\begin{warnbox}[101evler için önemli]
101evler tarafında \textbf{EUR (603)} ve \textbf{GBP (604)} kodları, resmi dokümantasyonda belirtilmediği için bizim tarafımızda \textbf{sıralı olarak atanmıştır}. Bu kodların 101evler sisteminde doğru karşılığa sahip olduğunu lütfen teyit ediniz; gerekirse doğru kodları tarafımıza iletin, hızla güncelleyelim.
\end{warnbox}

% ============================================================
\section{Filtreleme ve Atlama Kuralları}

Site veritabanındaki her ilan otomatik olarak feed'e dahil edilmez. Yalnızca aşağıdaki koşulları sağlayan ilanlar XML çıktısında yer alır:

\subsection{Ortak ön koşullar (her iki portal)}

\begin{enumerate}[leftmargin=*]
\item Yayın durumu \texttt{PUBLISHED} (Yayında) olmalı.
\item İlan editöründe ilgili portalın \textbf{export} (gönderim) bayrağı açık olmalı:
  \begin{itemize}
  \item 101evler için \texttt{export\_to\_101evler = true}
  \item hangiev için \texttt{export\_to\_hangiev = true}
  \end{itemize}
\item İlanın türü \textbf{Satılık} veya \textbf{Kiralık} olmalı (proje/günlük dahil).
\item İlanın \textbf{fiyatı} ve \textbf{para birimi} tanımlı olmalı (fiyat $> 0$).
\item Para birimi desteklenen değerlerden biri olmalı (TRY, USD, EUR, GBP).
\end{enumerate}

\subsection{101evler için ek koşullar}

101evler taksonomisinin kendine özgü olması nedeniyle aşağıdaki eşleştirmeler de yapılmış olmalıdır:

\begin{itemize}[leftmargin=*]
\item \textbf{type\_id} — 101evler emlak tipi kodu (Daire, Villa vb.)
\item \textbf{area\_id} — 101evler bölge kodu (şehir / mahalle)
\end{itemize}

Bu eşleştirmeler ilan editöründe ya da Sabit Listeler (\texttt{lookups}) panelinde 101evler tarafından sağlanan kod tablolarına göre yapılmaktadır.

\subsection{Atlama sebepleri}

Feed üreticisi her ilan için aşağıdaki kontrolleri sırayla uygular; başarısız olan ilan XML'e \textbf{eklenmez} ve dahili olarak sebep işaretlenir:

\begin{tabular}{ll}
\textbf{Kod} & \textbf{Anlamı} \\
\texttt{missing\_type\_id} & 101evler emlak tipi eşleştirilmemiş \\
\texttt{missing\_area\_id} & 101evler bölge kodu eşleştirilmemiş \\
\texttt{missing\_sale\_or\_rent} & Satılık/Kiralık bilgisi yok \\
\texttt{unsupported\_currency} & Para birimi desteklenmiyor \\
\texttt{missing\_price} & Fiyat tanımsız veya $\leq 0$ \\
\end{tabular}

% ============================================================
\section{Hata Ayıklama / Test Modu}

Her iki URL, \texttt{\&debug=1} parametresi eklendiğinde XML yerine \textbf{JSON özet} döner. Bu sayede portföyün hangi ilanlarının atlandığı görülebilir:

\begin{lstlisting}[style=xmlstyle]
GET https://alfa-emlak-websitesi.onrender.com/api/feeds/101evler.xml?token=<TOKEN>&debug=1

{
  "total": 142,
  "included": 138,
  "skipped": [
    { "listingId": "AE-2026-11001", "reason": "missing_area_id" },
    { "listingId": "AE-2026-11042", "reason": "unsupported_currency" }
  ]
}
\end{lstlisting}

Bu uç nokta entegrasyon sırasında doğrulama amaçlı kullanılabilir; ancak üretim ortamında düzenli olarak \textbf{çağrılmamalıdır}.

% ============================================================
\section{Önbellek (Cache) Davranışı}

Sunucu, XML çıktısını HTTP başlığı ile \texttt{Cache-Control: public, max-age=300, s-maxage=300} olarak servis eder. Yani:

\begin{itemize}[leftmargin=*]
\item Aynı URL'ye 5 dakika içinde yapılan tüm istekler aynı yanıtı döndürür.
\item Yayın değişikliklerinin (yeni ilan, fiyat güncellemesi, yayından kaldırma) feed'e yansıması \textbf{en fazla 5 dakika} sürer.
\item Çekim sıklığını 15--30 dakikaya ayarlamanız hem sizin hem de kaynak kullanımı açısından idealdir.
\end{itemize}

% ============================================================
\section{Görsel ve Medya}

\begin{itemize}[leftmargin=*]
\item Tüm görseller \textbf{HTTPS} adresleri üzerinden sunulur.
\item Görsel sıralaması \texttt{sort\_order} alanı ile belirlenir; en küçük değer kapak görselidir.
\item Görseller doğrudan portal tarafına \textbf{indirilebilir} (CDN üzerinden anonim erişime açık).
\item Sanal tur (\texttt{t\_360}) ve video (\texttt{youtube}) alanları, ilanda mevcutsa boş olmayan değerle gönderilir.
\end{itemize}

% ============================================================
\section{Sıkça Sorulan Sorular}

\subsection*{Bir ilanı portala manuel olarak ekleyebilir miyiz?}
Hayır, manuel ekleme gerekmiyor. Feed otomatik tüm yayında olan ilanları içerir.

\subsection*{Belirli ilanları feed'e dahil etmek istemiyorum, mümkün mü?}
Evet. Alfa Emlak panelinde her ilanın editöründe 101evler ve hangiev için ayrı ayrı gönderim bayrakları vardır. Bu bayrak kapatılan ilanlar feed'e girmez.

\subsection*{Çoklu dil desteği nasıl çalışıyor?}
Başlık (\texttt{ad\_title\_*}) ve açıklama (\texttt{ad\_description\_*}) alanları TR, EN, RU, DE ve FA dilleri için ayrı etiketlerle gönderilir. İlanın o dilde çevirisi yoksa ilgili etiket boş kalır veya hiç üretilmez.

\subsection*{İlan silindiğinde portala nasıl haber verilir?}
Silinmiş ya da yayından kaldırılmış ilanlar bir sonraki feed çekiminde XML'de \textbf{bulunmazlar}. Bunu portal tarafının fark edip ilanı pasifleştirmesi beklenir.

\subsection*{Aynı ilan iki kere geliyor mu?}
Hayır. \texttt{ad\_key} / \texttt{listing\_id} alanları benzersizdir; aynı ilan tek kayıt olarak gelir.

% ============================================================
\section{İletişim ve Destek}

Entegrasyon süreciyle ilgili her türlü teknik soru için iletişime geçebilirsiniz:

\begin{itemize}[leftmargin=*]
\item \textbf{Telefon:} +90 533 860 65 35
\item \textbf{Web:} \url{https://alfa-emlak-websitesi.onrender.com/tr/iletisim}
\item \textbf{Tasarım \& Yazılım:} Emy Software Studios — \url{https://emysoft.com}
\end{itemize}

\vspace{0.4cm}
\hrule
\vspace{0.4cm}
\begin{center}
\small\itshape\color{black!60}
Bu doküman, 101evler ve hangiev entegrasyonu için hazırlanmıştır. Belge versiyonu: 1.0 — \today
\end{center}

\end{document}
